% Generated draft: requires mathematical and source review.
\documentclass[11pt]{article}
\usepackage[margin=25mm]{geometry}
\usepackage{fontspec}
\setmainfont{Latin Modern Roman}
\usepackage{amsmath,amssymb,mathtools,unicode-math}
\setmathfont{Latin Modern Math}
\usepackage{enumitem,needspace,xurl,hyperref}
\setlength{\parindent}{0pt}
\setlength{\parskip}{5pt}
\setlength{\emergencystretch}{3em}
\newcommand{\N}{\mathbb{N}}
\newcommand{\Z}{\mathbb{Z}}
\newcommand{\Q}{\mathbb{Q}}
\newcommand{\R}{\mathbb{R}}
\newcommand{\C}{\mathbb{C}}
\newcommand{\F}{\mathbb{F}}
\newcommand{\A}{\mathbb{A}}
\newcommand{\PP}{\mathbb{P}}
\newcommand{\E}{\mathbb{E}}
\newcommand{\eps}{\varepsilon}
\newcommand{\dd}{\,\mathrm{d}}
\begin{document}
\section*{UnsolvedMath ID: 30006591}
\textbf{Source problem code:} catalogue entry
\subsection{Definitions and mathematical statement}
Let \(k\) be an integer for which the space \(S_k(\operatorname{SL}_2(\Z))\) of holomorphic cusp forms of weight \(k\) for the full modular group contains a normalized Hecke eigenform. Thus \(f\in S_k(\operatorname{SL}_2(\Z))\) is a simultaneous eigenvector for the Hecke operators and has Fourier expansion
\[
f(z)=\sum_{n\geq 1}a_f(n)e^{2\pi i n z},\qquad a_f(1)=1.
\]
The existence of such an \(f\) entails that \(k\) is even and that the cusp-form space is nonzero.

The supplied record does not define the normalization of \(\Lambda(f;s)\). Under the conventional completed modular \(L\)-function interpretation, set
\[
L(f,s)=\sum_{n\geq 1}\frac{a_f(n)}{n^s}
\quad\text{for sufficiently large }\operatorname{Re}(s),
\qquad
\Lambda(f;s)=(2\pi)^{-s}\Gamma(s)L(f,s),
\]
using analytic continuation to evaluate the completed function at the critical integers \(1,\ldots,k-1\). This normalization is an explicit interpretation requiring confirmation against the source.

Define the following \(\Q\)-vector subspaces of \(\C\):
\[
\begin{aligned}
L_f^+&=\operatorname{span}_{\Q}\{\Lambda(f;s):s\in\Z,\ 1\leq s\leq k-1,\ s\equiv 1\pmod 2\},\\
L_f^-&=\operatorname{span}_{\Q}\{\Lambda(f;s):s\in\Z,\ 1\leq s\leq k-1,\ s\equiv 0\pmod 2\}.
\end{aligned}
\]
In particular, the endpoints in the two lists are \(1,3,\ldots,k-1\) and \(2,4,\ldots,k-2\), respectively.

\textbf{Conjecture 1 (as supplied).} For every such weight \(k\) and every normalized Hecke eigenform \(f\in S_k(\operatorname{SL}_2(\Z))\), prove both equalities
\[
\dim_{\Q}L_f^+=\dim_{\C}S_k(\operatorname{SL}_2(\Z)),
\qquad
\dim_{\Q}L_f^-=\dim_{\C}S_k(\operatorname{SL}_2(\Z)).
\]

The supplied excerpt reports that functional relations prove the upper bounds
\[
\dim_{\Q}L_f^\pm\leq\dim_{\C}S_k(\operatorname{SL}_2(\Z)).
\]
This is a report of the source's claim, not an independent verification. The supplied curation describes equality as a folklore conjecture tied to the unproved Maeda conjecture; it supplies no precise formulation of that connection.

For background, the excerpt also uses the even and odd parts of a period polynomial \(P_f(X,Y)\):
\[
P_f^+(X,Y)=\frac12\bigl(P_f(X,Y)+P_f(-X,Y)\bigr),\qquad
P_f^-(X,Y)=\frac12\bigl(P_f(X,Y)-P_f(-X,Y)\bigr).
\]
It describes these as generating polynomials for critical values at odd and even integers, respectively. The supplied record does not define or normalize \(P_f\), so this background does not specify an additional mathematical assertion to prove.
\subsection{Short English statement}
For every normalized Hecke eigenform of full modular level, the conjecture says that its odd critical completed \(L\)-values and its even critical completed \(L\)-values each span a rational vector space whose dimension equals the complex dimension of the cusp-form space of the same weight. The completion convention requires source confirmation. The supplied record labels the problem open and reports an upper bound; these are supplied assessments, not independently verified conclusions.
\subsection{Sources}
\begin{itemize}
\item Supplied local UnsolvedMath registry record, permanent ID \(30006591\), title: \emph{Rational Dimensions of Critical Modular \(L\)-Values}; exact source problem code: \texttt{catalogue entry}. The record belongs to the Oberwolfach Reports Open Problems collection and gives proposed year 2026. No individual proposer is identified.
\item Supplied report citation: \emph{Multiple Zeta Values and Modular Forms} (2026), Oberwolfach Reports, DOI 10.4171/owr/2026/7, pp.~433--486. Source URL and supplied literature-source link, labelled \emph{Oberwolfach Report 7/2026}: \url{https://doi.org/10.4171/owr/2026/7}. The excerpt labels the assertion Conjecture~1; no precise page within the report is supplied.
\item Literature assessment dated \texttt{2026-08-22}: the supplied curation labels the problem open, reports only the upper bound as proved in the report, and describes equality as a folklore conjecture tied to the unproved Maeda conjecture. This is supplied dated curation, not independently verified literature assessment. The background explicitly requests independent verification of time-sensitive status and cited claims.
\item The registry supplies statement status \texttt{corrected\_verified} and states that the formulation was reconstructed and checked against the official report. That verification claim is part of the supplied record; no source visit or independent verification is claimed here.
\end{itemize}
\end{document}
